\section{Related Work}

The QBS framework sits at the intersection of several established literatures. Everettian decision-theoretic work derives or interprets Born-rule weights under rationality assumptions, including Wallace's formal decision-theoretic proof and Greaves's caring-measure interpretation. Work by Sebens and Carroll develops self-locating uncertainty as a basis for Everettian credence, while Saunders develops accounts of chance and probability in branching theories.

Quantum-suicide discussions provide historical background for extreme observer-selection intuitions, but they do not by themselves establish certain survival or quantum immortality. This manuscript therefore treats such thought experiments as background rather than as a theorem or implementation target.

Observer-selection proposals such as Hanson's mangled-worlds mechanism are closer to a physical selection model, but their mechanism differs sharply from the abstract accessibility map used here. Anthropically motivated decision theories, including Armstrong's anthropic decision theory, provide another nearby literature in which policy choice is studied under self-locating structure.

The provisional novelty claim of QBS is narrower than any claim to have introduced observer selection or Everettian decision theory. The distinctive combination is recognition-dependent policy change, policy-dependent trajectory and accessibility maps, exact decomposition of ordinary trajectory effects from first-person conditioning, and a separate analysis of cross-branch policy coherence. This positioning remains subject to revision as the literature review expands.
